\documentclass[conference]{IEEEtran}

\usepackage{amsmath,amssymb,amsfonts}
\usepackage{booktabs}
\usepackage{multirow}
\usepackage{hyperref}
\usepackage{cleveref}

\newcommand{\R}{\mathbb{R}}
\newcommand{\norm}[1]{\left\|#1\right\|}
\newcommand{\zstar}{z^{*}}
\newcommand{\Ahat}{\hat{A}}
\newcommand{\ecrit}{\varepsilon_{\mathrm{crit}}}

\begin{document}

\title{AEGIS: Supplementary Material}
\maketitle

\appendix

%% ===================================================================
\section{Proof Details}
\label{app:proofs}
%% ===================================================================

\subsection{Proposition 1: Optimal First-Order Structural Attack}

The first-order approximation gives $\Delta \zstar \approx S \cdot \mathrm{vec}(\delta A)$ where $S = (I - J_z)^{-1} J_A$. The optimization $\max_{\norm{\mathrm{vec}(\delta A)} \leq \varepsilon} \norm{S \cdot \mathrm{vec}(\delta A)}$ is a standard SVD problem: the maximum is $\varepsilon \cdot \sigma_1(S)$, attained by $\mathrm{vec}(\delta A^*) = \varepsilon \cdot v_1$ where $v_1$ is the leading right singular vector of $S$.

For the constrained case (symmetric, edge-only), we construct $S_c \in \R^{D \times |E|}$ with columns $S_{:,iN+j} + S_{:,jN+i}$ for each edge $(i,j)$. The constrained maximum is $\varepsilon \cdot \sigma_1(S_c)$. $\square$

\subsection{Proposition 2: Per-Node Robust Radius}

By Theorem~1(a), $\Delta \zstar_v \approx S_v \cdot \mathrm{vec}(\delta A)$ where $S_v$ is the block-rows of $S$ for node $v$. The change in the logit for class $y_v$ at node $v$ is:
\begin{equation}
|\Delta f_{y_v}| \leq \norm{\frac{\partial f}{\partial \zstar_v}} \cdot \norm{S_v} \cdot \norm{\delta A}_F.
\end{equation}
Misclassification requires $|\Delta f_{y_v}| \geq m_v$. Solving: $\norm{\delta A}_F \geq m_v / (\norm{\partial f / \partial \zstar_v} \cdot \norm{S_v}) = r_v$. Note that $S_v$ already contains the $(I - J_z)^{-1}$ factor, so no separate $(1-\rho)$ term appears. $\square$

\subsection{Non-Normality Index}

For a normal matrix $J_z$, $\norm{(I-J_z)^{-1}}_2 = 1/(1-\kappa)$ exactly (since $\kappa = \rho$ for normal matrices). For non-normal $J_z$, the resolvent norm can exceed the spectral-radius-based estimate by a factor $\eta = \norm{(I-J_z)^{-1}}_2 \cdot (1-\rho) \geq 1$. When $\eta \gg 1$, the operator exhibits ``transient amplification'': perturbations are temporarily amplified before contraction damps them~\cite{trefethen2005spectra}. Across our 9 datasets, $\eta \in [1.0, 1.4]$, indicating mild non-normality. The operator norm $\kappa = \norm{J_z}_2$ is therefore a faithful predictor of adversarial vulnerability in our setting, and is used in all tables and bounds.

%% ===================================================================
\section{Extended Experimental Results}
\label{app:extended}
%% ===================================================================

\subsection{Per-Seed Breakdown: Citeseer with Early Stopping}

Early stopping (best validation accuracy over 200 epochs, checked every 10 epochs) reduces Citeseer cert\% variance from $19.3\%$ to $4.0\%$:

\begin{table}[h]
\centering
\caption{Citeseer early stopping ablation (10 seeds). Accuracy values below reflect the corrected data loader; the main text (\cref{tab:cross_domain}) reports $66.0 \pm 0.7\%$ with early stopping on the same corrected loader.}
\small
\begin{tabular}{@{}lcc@{}}
\toprule
Metric & Without ES & With ES \\
\midrule
Accuracy      & $0.670 \pm 0.013$ & $0.660 \pm 0.007$ \\
Cert\%        & $72.0 \pm 19.3\%$ & $82.8 \pm 4.0\%$ \\
\bottomrule
\end{tabular}
\end{table}

\subsection{Mettack Comparison at Budget $k = 1, \ldots, 5$}

\cref{tab:mettack_full} reports IFT vs.\ Mettack damage at all budget levels on the corrected (BFS subgraph, sign-corrected gradient) implementation.

\begin{table}[h]
\centering
\caption{IFT vs.\ Mettack at budget $k=1,\ldots,5$ (Cora, single representative seed; 10-seed averages in the main text). IFT wins at all budgets.}
\small
\begin{tabular}{@{}rcccc@{}}
\toprule
$k$ & IFT & Mettack & Random & Brute-force \\
\midrule
1 & 0.355 & 0.052 & 0.078 & 0.355 \\
2 & 0.482 & 0.069 & 0.107 & 0.482 \\
3 & 0.586 & 0.079 & 0.117 & 0.586 \\
4 & 0.606 & 0.094 & 0.181 & 0.627 \\
5 & 0.616 & 0.100 & 0.192 & 0.661 \\
\bottomrule
\end{tabular}
\label{tab:mettack_full}
\end{table}

\subsection{Phase Transition Scan Data}

Full phase transition data (Cora, 50-node subgraph):

\begin{table}[h]
\centering
\caption{Phase transition: fixed-point shift vs.\ spectral radius.}
\small
\begin{tabular}{@{}cccc@{}}
\toprule
$\rho_{\text{target}}$ & $\rho_{\text{actual}}$ & Shift & Converged \\
\midrule
0.30 & 0.296 & 0.039 & Yes \\
0.50 & 0.495 & 0.044 & Yes \\
0.70 & 0.705 & 0.082 & Yes \\
0.85 & 0.854 & 1.255 & No \\
0.90 & 0.884 & 0.900 & No \\
0.95 & 0.936 & 2.018 & No \\
0.99 & 0.975 & 16.571 & No \\
\bottomrule
\end{tabular}
\end{table}

\subsection{Constrained vs.\ Unconstrained Tightness}

The gap between constrained ($\approx 1.00$) and unconstrained ($0.31$--$0.51$) tightness arises because the unconstrained SVD direction includes perturbations to non-edges, asymmetric entries, and self-loops---none of which are meaningful for undirected graph perturbation. The constrained sensitivity matrix $S_c$ restricts to symmetric edge-only perturbations, eliminating this artifact.

\begin{table}[h]
\centering
\caption{Constrained vs.\ unconstrained tightness (Cora, varying $N$).}
\small
\begin{tabular}{@{}rccc@{}}
\toprule
$N$ & Constrained & Unconstrained & Ratio \\
\midrule
20  & 1.009 & 0.310 & 3.3$\times$ \\
50  & 1.010 & 0.311 & 3.2$\times$ \\
100 & 1.014 & --- & --- \\
200 & 1.022 & --- & --- \\
\bottomrule
\end{tabular}
\end{table}

%% ===================================================================
\section{Implementation Details}
\label{app:implementation}
%% ===================================================================

\textbf{IGNN architecture.} Input projection $U \in \R^{d_{\text{in}} \times 64}$, state propagation $W \in \R^{64 \times 64}$ with spectral normalization, readout $h \in \R^{64 \times C}$. Operator: $F(Z) = \mathrm{ReLU}(\Ahat Z W^\top + U(X))$. Fixed-point iteration: max 50 steps, tolerance $10^{-5}$.

\textbf{ContractiveGCN-PF.} Same architecture as IGNN with 5 input features (is\_slack, is\_pv, P, Q, $|V|$) and 2-head readout ($\Delta|V|$, $\Delta\theta$). Binary adjacency from admittance matrix $Y$.

\textbf{Training.} Adam optimizer, lr=$0.01$ (graph benchmarks) / $10^{-3}$ (PF), weight decay $5 \times 10^{-4}$. Early stopping on validation accuracy over 200 epochs (graph benchmarks) / 30--50 epochs (PF).

\textbf{Subgraph extraction.} BFS from highest-degree node, capped at 50 nodes. The BFS approach guarantees connectivity (naive neighbor-slice yields near-empty subgraphs on high-degree hubs, e.g., WikiCS center has 3,324 neighbors).

\textbf{Randomized smoothing baseline.} Gaussian noise $\sigma = 0.01$ on adjacency entries, 200 Monte Carlo samples, Clopper-Pearson lower bound with $\alpha = 0.001$. Certified radius: $r_v^{\text{smooth}} = \sigma \cdot \Phi^{-1}(p_{\text{lower}})$.

\textbf{Mettack baseline.} Meta-Self variant: 2-layer GCN surrogate trained on IGNN pseudo-labels (100 epochs), gradient $-\partial \mathcal{L}_{\text{CE}} / \partial A$ for edge removal scoring. Corrected sign: negative gradient identifies edges whose removal increases misclassification.

\bibliographystyle{IEEEtran}
\bibliography{aegis}

\end{document}
